\section{GDF: auxiliary metrics and nuclear response}
\label{sec:gdf}

\subsection{Fitting AO-pair densities}

Gaussian density fitting represents AO-pair densities in an atom-centered
auxiliary basis $\{P\}$.  For a k-point pair $(\vk_i,\vk_j)$ with momentum
transfer $\vq=\vk_j-\vk_i$, the three-center source and two-center Coulomb
metric are
\begin{align}
  B_{P,\mu\nu}^{ij}
  &=\left(P\vq\middle|\mu\vk_i,\nu\vk_j\right),
  \label{eq:gdf-three-center}\\
  M_{PQ}^{\vq}
  &=\left(P\vq\middle|Q\vq\right).
  \label{eq:gdf-two-center}
\end{align}
With the metric factorization
\begin{equation}
  M^{\vq}=L^{\vq}L^{\vq\dagger},
  \label{eq:gdf-cholesky}
\end{equation}
the orthonormalized CDERI is
\begin{equation}
  C^{ij}=L^{-1}B^{ij}.
  \label{eq:gdf-cderi}
\end{equation}
A four-center integral is approximated by contracting two such CDERI factors.

\subsection{The three-center source derivative}

$B_{P,\mu\nu}$ contains an auxiliary center, a bra-AO center, and a ket-AO
center.  Its derivative for atom $A$ is
\begin{align}
  B_{P,\mu\nu}^{ij,Ax}
  ={}&\delta_{A,A_P}\partial_{P_x}B_{P,\mu\nu}^{ij}
    +\delta_{A,A_\mu}\partial_{\mu_x}B_{P,\mu\nu}^{ij}
  \notag\\
   &+\delta_{A,A_\nu}\partial_{\nu_x}B_{P,\mu\nu}^{ij}.
  \label{eq:gdf-source-centers}
\end{align}
If an integral library's \texttt{ip1}/\texttt{ip2} operator differentiates an
electronic coordinate, a translational identity may reverse the sign for a
nuclear-center displacement.  Each axis and atom slice must therefore be
checked with a small component test.

A periodic GDF builder may combine short-range real-space, long-range
reciprocal-space, and $\vq=0$ correction terms,
\begin{equation}
  B^{ij}=B_{\mathrm{SR}}^{ij}+B_{\mathrm{LR}}^{ij}
  +\delta_{\vq,0}B_{vbar}^{ij}.
  \label{eq:gdf-source-parts}
\end{equation}
The derivative must use the same screening, lattice images, and reciprocal
mesh as the central builder.

\subsection{Metric and Cholesky derivatives}

The two-center metric derivative contains the responses of both auxiliary
centers,
\begin{equation}
  M_{PQ}^{\vq,Ax}
  =\delta_{A,A_P}\partial_{P_x}M_{PQ}^{\vq}
   +\delta_{A,A_Q}\partial_{Q_x}M_{PQ}^{\vq}.
  \label{eq:gdf-metric-centers}
\end{equation}
Differentiating Eq.~\eqref{eq:gdf-cholesky} gives
\begin{equation}
  M^{Ax}=L^{Ax}L^\dagger+LL^{Ax\dagger}.
  \label{eq:gdf-cholesky-derivative}
\end{equation}
Define $W=L^{-1}M^{Ax}L^{-\dagger}$.  In the lower-Cholesky gauge,
\begin{equation}
  L^{Ax}=L\,\Phi(W),
  \label{eq:gdf-cholesky-gauge}
\end{equation}
where $\Phi$ retains the strict lower triangle and half of the diagonal.  The
CDERI derivative is therefore
\begin{equation}
  C^{Ax}
  =L^{-1}B^{Ax}-L^{-1}L^{Ax}C.
  \label{eq:gdf-cderi-derivative}
\end{equation}
The second term is the metric/auxiliary response.  Omitting it differentiates a
functional form different from the one used for the fitted central energy.

\subsection{J and K derivatives}

Using the diagonal k-pair CDERI, define the Coulomb auxiliary density as
\begin{equation}
  j_P=\frac{1}{N_k}\sum_{\vk\mu\nu}
      C_{P,\mu\nu}^{\vk\vk}D_{\nu\mu}^{\vk}.
  \label{eq:gdf-coulomb-aux-density}
\end{equation}
At fixed density,
\begin{equation}
  j_P^{Ax}=\frac{1}{N_k}\sum_{\vk\mu\nu}
      C_{P,\mu\nu}^{\vk\vk,Ax}D_{\nu\mu}^{\vk}.
  \label{eq:gdf-coulomb-aux-derivative}
\end{equation}
Consequently,
\begin{equation}
  J_{\mu\nu}^{\vk,Ax}
  =\sum_P\left(j_P^{Ax}C_{P,\mu\nu}^{\vk\vk}
              +j_PC_{P,\mu\nu}^{\vk\vk,Ax}\right).
  \label{eq:gdf-j-derivative}
\end{equation}
Exchange similarly applies the product rule to each of its two CDERI factors.
Reverse blocks of canonical k-pairs are reconstructed with the required
complex conjugation and AO transpose.

\subsection{Two appearances of auxiliary response in the gradient}

First, $C^{Ax}$ appears in the explicit target-Hartree derivative of
Eq.~\eqref{eq:hartree-gradient}.  Second, GDF J/K derivatives enter the source
HF CPHF source and Hessian action.  Adding metric response only to the explicit
target term while omitting it from the source response, or vice versa, breaks
agreement with the full-energy finite difference.

\subsection{The value-parity guard}

Define the difference between the central CDERI reconstructed by the response
builder and the CDERI stored by the actual SCF provider as
\begin{equation}
  \Delta_C=\max|C_{\mathrm{response\ builder}}-C_{\mathrm{SCF\ storage}}|.
  \label{eq:gdf-cderi-parity}
\end{equation}
If $\Delta_C$ exceeds tolerance, even plausible $B^{Ax}$ and $M^{Ax}$ are not
derivatives of the energy representation that was actually used, and the
gradient must be rejected.  Common causes are different values of $\omega$,
reciprocal meshes, auxiliary-basis ordering, k-pair orientation, or metric
factorization gauges.

\subsection{CPU/GPU boundaries}

The CPU package generates central and derivative primitives with the PySCF
periodic integral engine and performs linear algebra with NumPy/SciPy.  The GPU
package declares periodic two- and three-center primitive generation as a host
boundary, while CDERI storage, Cholesky/metric response, $dJ/dR$, $dK/dR$,
AX/Z-vector operations, and final contractions run on the CuPy device.
\texttt{device="gpu"} does not claim that integral metadata are also generated
on the GPU; it guarantees that supported numerical contractions do not
silently call the CPU implementation.

GDF completion requires full-energy finite differences, the translation sum,
and CDERI parity for distorted H$_2$O as well as H$_2$, for Gamma and
non-Gamma k-meshes, and at both SZV and DZVP levels.
